\chapter{Logs, and where everything is written}\label{ch:logs}

There are two logs, and confusing them is the most common way to look for a problem in the wrong
place.

\section{The application log}

One file per launch of ATE, in \texttt{log.directory}, named after the launch's start time:

\begin{lstlisting}[language=ATEShell]
/var/lib/ate/logs/ate-app-2026-11-01T09-02-31.log
\end{lstlisting}

It is opened in append mode and never rotated and never deleted, so each launch keeps its own
history and a Test Run is never split across files. One line per entry:

\begin{lstlisting}[language=ATEShell]
2026-11-01T09:02:31 INFO  ate: ATE 1.2.3 launched with config /etc/ate/config.json
\end{lstlisting}

Entries below \texttt{log.level} are dropped before they reach the file. The levels, in order, are
\texttt{debug}, \texttt{info}, \texttt{warning}, \texttt{error}; a level name the code does not
recognise is treated as \texttt{info}. What lands here is ATE's own activity: the launch, the
database refusal, interrupted runs closed, Device buffer overflows, the PDF sandbox decision, and
any exception that no code caught.

\section{The Run log}

Every line a Test Run produces --- what the schedule logged, what ATE logged about it, Device trace
--- is a row in the database's \texttt{log\_line} table, not a file. Each row has a level
(\texttt{DEBUG}, \texttt{INFO}, \texttt{WARNING}, \texttt{ERROR}), a source (\texttt{schedule},
\texttt{ate} or \texttt{device}), and an owner: the run as a whole, a Pre-test or Post-test Action,
or a step and Attempt. The Test Report prints only \texttt{INFO}, \texttt{WARNING} and \texttt{ERROR}
lines: \texttt{DEBUG} lines, and all Device trace, stay in the database. To see them, export the Run
log from History (CSV, JSON Lines or text) or query the database.

Device trace comes from \texttt{deviceTrace.enabled}, and each traced operation is one \texttt{DEBUG}
line, with the bytes shortened past \texttt{deviceTrace.maxBytes}. It is off by default for a
reason: a schedule that reads an instrument in a tight loop produces a line per read.

\section{When something inside ATE goes wrong}

PySide6 does not end the process when a slot or a thread's \texttt{run} raises: it hands the
exception to \texttt{sys.excepthook} and carries on, so a worker keeps recording while a screen
quietly stops updating. ATE installs a hook that writes such an exception to this launch's
application log, as \texttt{Unhandled exception on thread}, with the traceback, and tells the
Operator in the status bar:

\begin{lstlisting}[language=ATEShell]
Internal error: this window may have stopped updating. A Test Run in progress keeps recording. Details: /var/lib/ate/logs/ate-app-2026-11-01T09-02-31.log
\end{lstlisting}

If the log could not be written, the same line says \texttt{The details could not be written to the
application log} instead of naming a path. The status bar line appears only from the first error
onward and stays until ATE is closed, because the screen may be wrong from that point on. It is
selectable text, so the path can be copied.

What that hook does \emph{not} see: an exception leaving a plain \texttt{threading.Thread} goes to
\texttt{threading.excepthook} instead, and the Device reader threads --- the only such threads ---
catch their own. An exception inside a Test Run's worker is caught by the run engine first and
reported through that run's own log and Test Report. So a Test Run failure is found in the Run log,
and an application failure in the application log.

\section{Device warnings worth reading}

ATE reads every TCP, SERIAL and UDP Device in the background and buffers what arrives, up to
\texttt{devices.receiveBufferBytes}. When the buffer overflows it drops the oldest bytes, records one
WARNING in the Run log with source \texttt{device}, and the next read result carries
\texttt{overflow = true}:

\begin{lstlisting}[language=ATEShell]
TCP 127.0.0.1:5025: receive buffer overflow - dropped 4096 oldest byte(s)
\end{lstlisting}

A Test Run that ends with Devices still open closes them and records one WARNING in the Run log per
Device, naming it as \texttt{left open at end of run}. That is untidy, not dangerous. On an aborted
run the closes are given three seconds in total; a close the operating system holds longer than that
is left running on its own thread and logged as \texttt{was still closing when the aborted run ended}
--- it does close, just later (chapter~\ref{ch:troubleshooting}).

\cando{find the line that explains a launch refusal, and tell the application log apart from the Run
log when someone asks "where is the log of this run?".}
